% The spoken body of the homily for the feast of Our Lord Jesus Christ the
% King. Continuous preaching, ready to read aloud: no heading, no direction to
% the preacher, no citation inside the speech. The loci, the relation to the
% reviewed readings, the audience, the word count and the pace estimate stand
% in the terminal note, which no one reads out. The homily is preached after
% the Gospel, so the chants and prayers before it are spoken of as heard and
% those after it as still to come. The blank lines are the paragraph breaks a
% speaker recovers his place by, and the vertical spaces divide the speech
% into its movements: the throne and the prisoner, and the question; "not of
% this world" as we hear it; "not from hence", with Augustine and Aquinas;
% how the kingdom takes hold, with Chrysostom, the Epistle and the Collect;
% why the King is bound, with Chrysostom, the Offertory, Aquinas and
% Augustine; the altar, with the Secret and the Preface; what the King asks,
% with Augustine on the Gradual's kings, the Introit's verse and Bellarmine;
% the Postcommunion's soldiers, with Hilary and Schuster; the week's
% response; the hope, with the Alleluia, the Communion, Augustine and
% Schuster; and the return to the governor's hall and the throne. The
% honorific is tied to the name so that "St" never ends a line.
\raggedcolumns

``Art thou a king then?''

It is early morning, before the Pasch. The accusers will not enter the
governor's hall, so Pilate goes in to the prisoner, bound since his arrest in
the garden, and puts to him the question on which this whole feast turns.

Remember how this Mass began. The Introit took us up to the throne of heaven,
where the angels cry, ``The Lamb that was slain is worthy to receive power and
divinity and wisdom and strength and honour.'' The Gradual sang of a King who
``shall rule from sea to sea,'' whom ``all kings of the earth shall adore,''
whom ``all nations shall serve.'' And then the Gospel set the same King before
us bound, in a Roman governor's hall, saying, ``My kingdom is not of this
world.''

So which is it: a King who rules from sea to sea, or a kingdom not of this
world?

\medskip

Many of us have settled that question without noticing. We hear ``My kingdom
is not of this world,'' and take it to mean that Christ is King somewhere
else: in heaven, in the soul, in this church for an hour on Sunday. Then
Monday comes, with the office, the household and the ballot box, and we
suppose those belong to some other king, or to none.

But listen to what he actually says. St~Augustine noticed what is easy to
miss. The Lord did not say, my kingdom is not \emph{in} this world. He said it
is not \emph{of} this world. And at the end of the sentence he did not say, my
kingdom is not \emph{here}. He said, ``my kingdom is not \emph{from hence}.''
From hence. The words tell us where the kingdom comes from, not where it is.
``For His kingdom is here until the end of the world,'' Augustine says,
``having tares intermingled therewith until the harvest.'' It is here; it is
not from here. ``It only sojourns as a stranger in the world.''

Think what follows. A kingdom that comes from one place in the world is
bounded by that place; its borders are where its armies stop. A kingdom whose
origin is not in the world at all is bounded by no place. St~Thomas Aquinas
puts it simply: it is not from here, and yet it is here, because it is
everywhere. That is how a bound prisoner can claim the whole earth. There is
no Monday outside his kingdom.

\medskip

Then how does such a kingdom take hold of anyone? ``If my kingdom were of this
world, my servants would certainly strive that I should not be delivered to
the Jews.'' St~John Chrysostom drew the lesson: ``Here He shows the weakness
of kingship among us, that its strength lies in servants; but that which is
above is sufficient for itself, needing nothing.'' This kingdom gathers its
people another way: ``Every one that is of the truth heareth my voice.'' It is
entered by hearing. St~Paul, in the Epistle, gave thanks to the Father who
``hath translated us into the kingdom of the Son of his love.'' Translated:
carried over. We were not conquered into this kingdom. We were carried into
it.

The Collect asked that the nations be subjected to Christ, but hear how. The
nations are ``separated by the wound of sin'': what scatters mankind is first
of all a wound. And the rule they are to be brought under is ``His most loving
dominion.''

\medskip

Why, then, is the King bound? If his power is greater than Pilate's, how was
he made a captive? Chrysostom answers, ``By consenting, and giving Himself
up.''

Christ was born a King, as Chrysostom insists. Yet this Mass says that he
receives. In the Offertory, said as the gifts are offered, the Father says to
his Son, ``Ask of me, and I will give thee the Gentiles for thy inheritance.''
Why should the eternal Son ask? St~Thomas answers that God wills even what is
owed to be received by asking, and that Christ asked by his prayer and by his
passion. St~Augustine says that this asking belongs to the Son made man, ``who
offered up Himself as a Sacrifice in the stead of all sacrifices.'' The
kingdom is his by birth. As man, he wins the nations in that hall, by the
passion that has begun. The Lamb is worthy because he was slain.

That sacrifice is not behind us. In the Secret, the prayer said quietly over
the gifts, the priest will offer ``this host for the reconciliation of
humanity,'' and ask that Christ bestow on all the nations ``the gifts of unity
and peace'': unity for those whom sin has scattered. The wound the Collect
named is healed at this altar. Then the Preface will name the kingdom by what
it holds: ``a kingdom of truth and life; a kingdom of holiness and grace, a
kingdom of justice, love and peace.'' The first good it names is the truth the
prisoner was born to bear witness to.

\medskip

What, then, does this King ask of us?

Look again at the Gradual's kings. The psalm goes on to give the reason for
their homage: ``For he shall deliver the poor from the mighty.'' That poor
man, says St~Augustine, ``is the people of men believing in Him. In this
people are also kings adoring Him.'' The king who adores Christ does not stand
above the poor. He kneels among them and confesses his sins. That is our first
place in this kingdom.

Augustine gives the kings a second place: their authority, he says, leads men
into the fellowship of the Church. Most of us wear no crown, yet every one of
us rules something: a household, a classroom, a ward through the night, a
vote. Each is an office, and the King whose kingdom is not from hence claims
every one. The Introit's verse is a prayer fit for every such office: ``Give
to the king thy judgment, O God, and to the king's son thy justice.''
St~Robert Bellarmine, on that verse, says what authority is for: ``kings were
created for the people, not the people for kings.'' Christ's kingdom takes no
one's throne from him. It tells him for whom he holds it.

\medskip

After Communion we shall pray as those ``who glory to fight under the standard
of Christ the King'': soldiers, a few minutes after a Gospel in which the King
says that his servants do not fight for him. So what is this fight? It is not
his defence; he needs none. It is our share in his victory over sin, for
St~Hilary says that the kingdom of God is where sin is overcome and death
destroyed. Blessed Ildefonso Schuster, on this prayer, names the standard:
``The Royal Standard of Jesus Christ is the Cross.'' We fight with the King's
own weapons in Pilate's hall: the truth, patience, and a life given up.

So take one thing out of this church. This week, choose one place where you
hold some authority, however small. Before you use it, make the Introit's
verse your own: Give to the king thy judgment, O God. Ask for whom you hold
it. Ask whose voice you are hearing. ``Every one that is of the truth heareth
my voice.''

\medskip

We shall not see this kingdom complete. Some give him their homage willingly
now; all will give it at the judgment, when every knee bends. The Alleluia
sang that ``his power is an everlasting power that shall not be taken away,''
and the Communion will sing, ``the Lord shall sit king for ever \dots\ the
Lord will bless his people with peace'': not peace in this world, Augustine
says, but peace in himself. Schuster says of that chant that ``at the moment
of Holy Communion Jesus is enthroned in the faithful heart.'' The throne of
heaven and the throne in a faithful heart belong to one King.

\medskip

Come back, then, to the governor's hall. The prisoner's first words to Pilate
were a question: ``Sayest thou this thing of thyself, or have others told it
thee of me?'' He asks us the same. Do we call him King of ourselves, or only
because others have told us? Soon he will be on this altar, the King who gave
himself up. Be of the truth, and hear his voice. Give him your Monday as well
as your Sunday. His kingdom is not from hence, and so there is nowhere it does
not reach.

With the angels of the Introit, then, we say it of the prisoner of the Gospel:
``The Lamb that was slain is worthy to receive power and divinity and wisdom
and strength and honour.'' ``To him be glory and empire for ever and ever.''
